\documentclass[11pt,a4paper]{article}
\usepackage[margin=23mm]{geometry}
\usepackage[T1]{fontenc}
\usepackage{amsmath,amssymb,booktabs,hyperref,xurl}
\usepackage[numbers,sort&compress]{natbib}
\newcommand{\PrimaryPairs}{93594}
\newcommand{\IntrinsicFWHM}{0.336\%}
\newcommand{\ThousandENCResolution}{2.54\%}
\newcommand{\SpeedLow}{0.174}
\newcommand{\SpeedHigh}{0.261}
\newcommand{\LeakageShotOneSecond}{61.2}

\title{Assessment of the SeF$_6$ ion-TPC proposal:\\noise, collection time, pile-up and calorimetry}
\author{NAUSICAA working note}
\date{4 November 2026}
\begin{document}
\maketitle

\section{Assessment and scope}
Nygren et al.~\cite{Nygren_2018} establish an attractive physics case:
the high $^{82}$Se $Q$ value suppresses common gamma backgrounds, thermal
ions diffuse little, and direct collection avoids multiplication noise.
The difficult part is simultaneously obtaining fine tracking, complete
charge collection and low-noise calorimetry for slow, distributed charge.
The original sub-50-electron ENC is an engineering projection. Later
Topmetal-S measurements in the local literature are about 115--140
electrons, and the N$\nu$DEx conceptual design explicitly identifies
two-neutrino coincidence backgrounds for enriched gas.

This note distinguishes measured electronics performance, design assumptions,
and analytical estimates. It is not a new SeF$_6$ transport or track simulation.
The calculations use $Q=2.995$ MeV, $W=32$ eV and $F=0.19$ from the original
paper; $W$ and $F$ were inferred using SF$_6$, rather than measured in the
proposed high-pressure SeF$_6$ detector. Track extents below are illustrative.

\section{What supports the CSA noise assumption?}
Section 4 and figure 7 of Nygren et al. propose a collection electrode smaller
than approximately 1 mm, a few pF of input capacitance, a femtofarad feedback
capacitor and an almost-off feedback transistor with adjustable resistance.
Charge focusing allows a small collection electrode at a several-mm to cm
pixel pitch. Direct waveform digitization replaces a long analog shaper;
timing comes from the rising edge and energy from digital filtering. They
require a long signal retention time and good noise at that integration time.
They do not give a measured ENC-versus-shaping-time curve for the proposed
few-pF electrode, a complete parallel/series noise budget, or a calibrated
MeV calorimetric demonstration at 10 bar.

Their cited Topmetal-II$^-$ demonstration~\cite{An2016TopmetalII} reported
13.9 electrons ENC with a 1 ms digital trapezoidal filter and 7.6 ms feedback
decay. This used a $25\times25~\mu$m collection node and approximately
23 fF input capacitance. The injection capacitance of 5.5 fF came from
simulation and was explicitly not independently calibrated. This is evidence
for direct ion sensing, but it does not establish the same ENC for a roughly
3 pF electrode and a much longer input pulse.

The later local papers provide more concrete information:
\begin{itemize}
\item Yang et al.~\cite{Yang2024TopmetalCSA}, sections 2--3, describe a
130 nm CMOS folded-cascode CSA, PMOS input transistor chosen for flicker
noise, adjustable feedback and external digitization. They measure input
capacitance near 3 pF, injection capacitance near 250 fF, gain around
220 mV/fC and ENC around 115 electrons after digital trapezoidal shaping.
They explicitly state that their required ENC below 45 electrons for the
target resolution has not yet been met.
\item Liang et al.~\cite{Liang2024TopmetalS}, section 2.2, measure
$121\pm9$ electrons ENC after calibrating the injection capacitance against
external capacitors. They report $255\pm10$ fF and explain that earlier
underestimation of this capacitance led to underestimated noise in earlier
Topmetal-S reports. This finding should not be applied as an unmeasured
numerical correction to Topmetal-II$^-$, which is a different device.
ENC rises to about 140 electrons as feedback decay is shortened from
200 to 10 ms; the pulse rise time is about 3 ms. The tested feedback decay
range extends to about 1.5 s. A long decay time alone does not establish
low ENC for every slow charge waveform. Ion demonstrations use air and
SF$_6$ near atmospheric pressure, not a complete 10-bar SeF$_6$ calorimeter.
\item The 2025 low-leakage study~\cite{Wang2025LowLeakage}, sections III--IV,
demonstrates charge retention of tens of seconds and sub-0.6 fA inferred
leakage. Its low-noise CSA plus switch performance is a separate simulation:
29 electrons for the CSA, 487 with the switch, and 39 after an ideal noiseless
CR--RC filter. The simulated filtered ENC grows to 52 electrons as the
injection duration reaches 300 ms. Thus 39 electrons is not a measured
system ENC for second-long high-pressure events.
\end{itemize}

The physical difficulty is the trade-off between bandwidth, input capacitance,
leakage/feedback noise, flicker noise and baseline stability. Suppressing white
series noise by longer filtering does not also suppress all parallel or
low-frequency noise. As an illustrative charge-counting limit only, independent
Poisson leakage contributes
\begin{equation}
 \sigma_{\rm leak}^2=I_{\rm leak}T/e.
\end{equation}
For a constant 0.6 fA leakage current integrated for 1 s this gives
\LeakageShotOneSecond\ electrons RMS. This is not a prediction for the
switch circuit: its feedback correlations and actual filter must be included.

\section{Track length and drift times}
Nygren section 5 compares 10-bar gases at their respective endpoints:
about 2.99 MeV for $^{82}$Se and 2.46 MeV for $^{136}$Xe. The paper argues
for less multiple scattering in the lower-$Z$ mixture and longer tracks at
the higher endpoint. It does not provide a quantitative same-energy range
comparison or distributions of track length and $\Delta z$ for both gases.
Figure 8 illustrates centimetre-scale extents; the CDR uses approximately
10 cm as a characteristic 0$\nu\beta\beta$ event size at 10 bar.

Consequently, the qualitative topology advantage is supported, but a generic
claim that SeF$_6$ has a longer electron stopping range at equal pressure and
equal energy is not established by these papers. Equal pressure fixes
molecule/atom number density at a fixed temperature, not mass density or
stopping power. Less multiple scattering, true path length, endpoint
separation and projection along $z$ are distinct quantities. A quantitative
gas comparison would require matched transport simulations and orientation
distributions. One displayed event is not an ensemble range measurement.

For an ion mobility $K$ and field $E$, distinguish full drift latency from
the arrival spread of one species across a track:
\begin{equation}
 v=KE,\qquad t_{\rm drift}=L/v,\qquad
 \Delta t_{\rm track}=\Delta z/v.
\end{equation}
For reduced mobility $K_0$ referenced to $(P_0,T_0)$,
\begin{equation}
 K=K_0\frac{P_0}{P}\frac{T}{T_0}.
\end{equation}
Using Nygren's approximate $K_0=0.4$--$0.6$ cm$^2$/(V s), 10 bar,
293.15 K and 400 V/cm gives \SpeedLow--\SpeedHigh\ m/s. This extrapolates
SF$_6$ information to SeF$_6$. The CDR adopts a speed above 20 cm/s and
approximately 7 s full drift over 1.6 m~\cite{NnDEx-100:2023alw}. The
original paper's 0.5--1 s ton-scale drift estimate and mm/ms language should
not be used without this density correction; they are optimistic compared
with its own reduced mobilities at 10 bar and the later CDR.

\begin{center}\small
\begin{tabular}{lrr}\toprule
Extent $\Delta z$ (mm) & $v=0.1$ m/s: time (s) & $v=0.2$ m/s: time (s)\\\midrule
10 & 0.1 & 0.05\\
20 & 0.2 & 0.1\\
100 & 1 & 0.5\\
200 & 2 & 1\\
1600 & 16 & 8\\
\bottomrule\end{tabular}

\end{center}
The 0.1 m/s column is the current NAUSICAA adopted Xe-ion speed, a design
assumption rather than a measured equality of Xe and SeF$_6$ mobilities.
Thus 10--20 mm implies 0.05--0.2 s at these representative speeds. A
1--2 s interval is appropriate for 10--20 cm at 0.1 m/s, or requires a
slower speed of 1 cm/s for the original 10--20 mm extent. Diffusion,
multiple ion species and electronics response can broaden the observed
waveform beyond the single-species geometrical estimate.

An additional primary-source check is the 2026 simulation framework
\cite{Liang2026NvDExSimulation}, which calculates reduced SeF$_5^-$ and
SeF$_6^-$ mobilities of 0.444 and 0.430 cm$^2$/(V s). These are calculations,
not high-pressure SeF$_6$ measurements. It explicitly treats the two species
as a hypothesis because their yields and the lifetime of SeF$_6^-$ in pure
high-pressure gas remain unverified.

\section{Pile-up and the absence of prompt timing}
The relevant event rate includes the entire low-energy background and
two-neutrino spectrum, not just the very small single-event rate at $Q$.
For a local Poisson rate $r$ and unresolved time interval $T$, the probability
of another event is $1-e^{-rT}$, approximately $rT$ at low occupancy.
The ambiguous spacetime volume also matters: simultaneous events on distant
channels need not be confused. Full drift latency does not equal mandatory
global acquisition dead time in a continuously digitized TPC.

The CDR section IV.E explicitly estimates the coincidence problem. Assuming
events are separable when approximately $10\times10\times10$ cm$^3$ apart,
its approximate $2\nu\beta\beta+2\nu\beta\beta$ background in the ROI is
0.06 events/year for 100 kg natural SeF$_6$ and about 8 events/year for
enriched gas. The latter exceeds its approximately 0.4 events/year
single-event gamma estimate before topological rejection. These are CDR
projections under their energy-window and separation assumptions, not new
rates inferred from our timing table. Their proposed remedies include
topology studies and potentially prompt scintillation readout, whose light
yield and practical implementation require study.

Two distinct negative-ion mobilities can in principle provide an absolute
drift distance from
\begin{equation}
 \Delta t_{12}=L(1/v_1-1/v_2).
\end{equation}
This helps fiducialization without a prompt $t_0$, but requires identifying
and pairing both species. It does not automatically associate all charge
unambiguously when two extended tracks overlap. The 2026 framework retains
this dual-species assumption explicitly. A complete pile-up assessment must
include the waveform, threshold, local occupancy and pairing ambiguity.

\section{Energy-resolution budget}
For $N=Q/W=\PrimaryPairs$ primary charges and one complete-event charge
estimate per selected channel, assume independent identical RMS noise ENC.
Adding independent variance gives
\begin{equation}
 R_{\rm FWHM}=2\sqrt{2\ln2}\,
 \frac{\sqrt{FN+n\,\mathrm{ENC}^2}}{N}.
 \label{eq:budget}
\end{equation}
This is an optimistic instrumental budget: it assumes complete collection,
perfect gain/baseline calibration, no thresholds, no recombination or
attachment fluctuations, no pile-up and independent channel estimates.
The intrinsic FWHM is \IntrinsicFWHM.

\begin{center}\small
\begin{tabular}{lrrr}\toprule
Independent channels $n$ & 30 electrons ENC & 50 electrons ENC & 120 electrons ENC\\\midrule
25 & 0.505\% & 0.713\% & 1.546\%\\
100 & 0.826\% & 1.302\% & 3.038\%\\
300 & 1.350\% & 2.205\% & 5.240\%\\
400 & 1.546\% & 2.538\% & 6.048\%\\
\bottomrule\end{tabular}

\end{center}

The CDR's 25 channels at 30 electrons reproduces its quoted 0.51\% FWHM.
Nygren's few-hundred-pixel assumption yields a considerably less favorable
budget even at 50 electrons. In particular, a total electronic noise near
1000 electrons gives \ThousandENCResolution\ FWHM including Fano noise;
being below 1000 electrons is not by itself sufficient to ensure below 2\%.
The maximum independent-channel budgets before additional imperfections are:

\begin{center}\small
\begin{tabular}{lrrr}\toprule
Target FWHM & Total noise budget (electrons) & Max. $n$, ENC 50 & Max. $n$, ENC 120\\\midrule
0.5\% & 147.3 & 8.7 & 1.5\\
0.7\% & 244.2 & 23.8 & 4.1\\
1.0\% & 374.4 & 56.1 & 9.7\\
2.0\% & 783.6 & 245.6 & 42.6\\
\bottomrule\end{tabular}

\end{center}
These are continuous variance budgets; an integer channel limit is their floor.

For correlated channel/time estimates, replace $n\,\mathrm{ENC}^2$ by
$\boldsymbol w^T C\boldsymbol w$, using the actual estimator weights and
noise covariance. Sample counts cannot be equated to independent noise
measurements. Conversely, a short injected-pulse ENC cannot simply be used
once per pixel for an extended, time-dependent ion current without verifying
the complete estimator. Threshold selection can discard low-density track
charge and produce energy bias and non-Gaussian tails; the original variance
budget does not cover these effects. Measuring both polarities can reduce
independent instrumental noise but not the shared Fano fluctuation.

\section{Implications and reproducibility}
The strongest demonstrated advantages are high endpoint energy and thermal
ion diffusion. The limiting engineering evidence concerns noise at realistic
input capacitance and waveform duration, reliable collection and charge
reconstruction, and enriched-gas pile-up. A convincing calorimetric validation
should inject known charge with realistic time distributions, measure the
full channel covariance and baseline stability, and report total reconstructed
energy as channel multiplicity and duration vary. A matched SeF$_6$/Xe track
study would resolve the separate question of $z$-extent distributions.

Numerical inputs are in \path{config/nygren_review.toml}. Julia saves SI-based
calculations, source hashes and code in \path{results/nygren_review/budgets.h5};
Python presents the saved values. Reproduce from the repository root:
\begin{verbatim}
JULIA_DEPOT_PATH="$PWD/work/julia-depot:$HOME/.julia" \
  julia --startup-file=no --project=. scripts/validation/nygren_review.jl
python3.13 scripts/validation/present_nygren_review.py
latexmk -pdf -interaction=nonstopmode -halt-on-error \
  -outdir=latex_generated/nygren_review docs/nygren_ion_tpc_review.tex
\end{verbatim}
\bibliographystyle{unsrtnat}
\bibliography{notes/physics/references}
\end{document}
